\SourceLeaf{317}%
The final values are
\[
\begin{aligned}
\log\lambda&=2.7315487\\
\log l&=2.7315438\\
l&=538''\,9442=0^0\,8'\,58''\,9442.
\end{aligned}
\]

\indent The full convenience of this procedure only becomes apparent once one has mastered the little aids that make up the working machinery of methods of this kind. Instructions on that would be out of place here. I will only point out that what looks like four calculations in the example above should be written as a single calculation, with only the last digits filled in or corrected at each new pass. In any case, one need keep only the last calculation. This is a great advantage, especially in surveys of considerable extent, since the entire essential part of the calculation for all the triangle sides is then available in the smallest possible space and in the clearest form for checking its accuracy whenever desired.

\Article{21}
I now turn to the main problem itself. It calls for a method for the surface of the ellipsoid like the one given above for the surface of the sphere. This rather intricate problem will be solved here in two entirely different ways. The first derivation is based on the conformal mapping from the ellipsoid onto the sphere. The theory of this mapping was developed in the first memoir. Finding this solution may therefore be regarded as the \textit{first indirect} use of that theory for higher geodesy. (See Art. 11.)

Let \(B+\frac12b\) and \(B-\frac12b\) now denote the latitudes of two points on the ellipsoid, with their difference in longitude denoted by \(l\), the segment of a geodesic line between them (measured here in any unit) by \(r\), and the azimuths of the line at the first and second endpoints by \(T+\frac12t\) and \(T-\frac12t\pm180^0\). The problem, then, is to find \(b,l\) and \(t\) from \(r,B\) and \(T\) by formulas analogous to those given above
